\documentclass[10pt,a4paper]{amsart}
\usepackage[margin=19mm]{geometry}
\usepackage{amsmath,amssymb,mathtools}
\usepackage[scr=rsfs,cal=euler]{mathalfa}
\usepackage{xfrac}
\usepackage[unicode,hidelinks,bookmarks=false]{hyperref}
\newcommand{\ie}{i.e.\ }
\newcommand{\cf}{cf.\ }
\newcommand{\resp}{respectively\ }
\newcommand{\Subsection}{\subsection}
\providecommand{\nobreakdash}{\nobreak\hbox{-}}
\setlength{\emergencystretch}{2em}
\multlinegap0pt
\usepackage{amssymb}
\newtheorem{theorem}{Theorem}
\title{Poisson–Hamiltonian Pontryagin Dynamics and Optimal Control of Mechanical Systems on Lie Groupoids}
\author{Ghorbanali Haghighatdoost}
\date{Source: arXiv:2602.20326v1 (CC BY 4.0)}
\begin{document}
\maketitle

\noindent\emph{Source and licence.} Poisson–Hamiltonian Pontryagin Dynamics and Optimal Control of Mechanical Systems on Lie Groupoids. Version arXiv:2602.20326v1 (CC BY 4.0), distributed by arXiv under the Creative Commons Attribution 4.0 International licence, http://creativecommons.org/licenses/by/4.0/, as recorded in the arXiv licence metadata for this exact version and shown on the abstract page https://arxiv.org/abs/2602.20326v1. Full licence notice: \texttt{licenses/arxiv-2602.20326-CC-BY-4.0.txt}.\par\smallskip

\noindent\textit{Editorial scope.} Counted unit: the article's theorem theorem (source0.tex lines 363--391). The article's complete original proof of it is delivered with it (source0.tex lines 392--448, one \texttt{proof} environment of 57 lines). No mathematics is written, repaired or re-derived: the statement and the proof are the article's own lines, in the article's own notation, and no retained line is substituted or re-flowed.  No further source-local context is needed: every cross-reference of every retained line resolves inside the two delivered spans.  Nothing was omitted from any retained span, so the package carries no omission record.  No retained line contains a citation, so the wrapper carries no bibliography and no bibliography entry.  No retained line contains a figure environment, an image-inclusion command or drawing code, so no image is delivered and the package declares no asset dependency.  Editorial formatting only, in addition: an AMS \texttt{amsart} preamble that restates the article's own declarations and macros used by the retained lines, namely the environments \texttt{theorem} on separate counters, together with the standard packages those lines need.  The retained environments print in this document as Theorem 1 in order of appearance, whereas the article prints the same items with the numbering of its own full document.  The complete native source of the article as distributed in the arXiv e-print for this exact version is bundled unchanged as \texttt{sources/195179/source0.tex}.
\begin{theorem}[Equivalence of variational and Hamiltonian optimal control]
Let $A \to M$ be a Lie algebroid with anchor $\rho : A \to TM$, and let
\[
\Phi : E \to A
\]
be a smooth fiberwise surjective bundle map defining a control system on $A$.  
Let $L : E \to \mathbb{R}$ be a smooth Lagrangian, regular in the control variables.

Then the following statements hold:

\begin{enumerate}
\item[(i)] Every critical trajectory of the variational optimal control problem on $A$, defined by minimizing
\[
J(u) = \int_0^T L(u(t))\,dt
\]
over admissible curves, projects to an integral curve of the Hamiltonian vector field generated by the reduced Pontryagin Hamiltonian
\[
\mathcal{H}(\alpha_x) = \langle \alpha_x, \Phi(u^*(\alpha_x)) \rangle - L(u^*(\alpha_x))
\]
on the dual Lie algebroid $A^*$, where $u^*$ is determined by the stationarity condition.

\item[(ii)] Conversely, every solution of the Hamiltonian system on $A^*$ generated by $\mathcal{H}$ projects to a critical trajectory of the variational optimal control problem on $A$.

\item[(iii)] The Hamiltonian flow generated by $\mathcal{H}$ is a Poisson flow with respect to the canonical linear Poisson structure on $A^*$, and its trajectories are confined to symplectic leaves of $A^*$.

\end{enumerate}

Consequently, the variational and Hamiltonian formulations of the optimal control problem on the Lie algebroid $A$ are equivalent.
\end{theorem}

\begin{proof}[Proof sketch]
We outline the standard argument.

\medskip
\noindent\textbf{Step 1 (Pontryagin bundle and stationarity).}
Consider the Pontryagin bundle $A^*\times_M E$ and the Pontryagin Hamiltonian
\[
H(\alpha_x,u_x)=\langle \alpha_x,\Phi(u_x)\rangle - L(u_x).
\]
Regularity of $L$ in the control variables implies that the stationarity condition
\[
\partial_u H(\alpha_x,u_x)=0
\]
determines (locally, and globally under appropriate convexity) a smooth feedback $u^*(\alpha_x)$. Hence the reduced Hamiltonian
\[
\mathcal H(\alpha_x)=H(\alpha_x,u^*(\alpha_x))
\]
is well-defined on $A^*$.

\medskip
\noindent\textbf{Step 2 (Variational problem $\Rightarrow$ Hamiltonian system).}
Let $u(t)$ be a critical control for the variational problem and set $a(t)=\Phi(u(t))$ with base curve $x(t)$. Introduce a Lagrange multiplier curve $\alpha(t)\in A^*_{x(t)}$ enforcing the admissibility constraint $\dot x=\rho(a)$ and consider the augmented functional
\[
\widetilde J(u,\alpha)=\int_0^T\Big(L(u(t))-\langle \alpha(t),\Phi(u(t))\rangle
+\langle \alpha(t),a(t)\rangle\Big)\,dt
=\int_0^T\big(-H(\alpha(t),u(t))+\langle \alpha(t),a(t)\rangle\big)\,dt.
\]
Taking first variations with respect to $u$ gives the stationarity condition
$\partial_u H(\alpha(t),u(t))=0$, hence $u(t)=u^*(\alpha(t))$.
Variations compatible with the Lie algebroid constraint yield the adjoint equations; equivalently, the pair $(x(t),\alpha(t))$ satisfies Hamilton's equations on the Poisson manifold $A^*$ generated by $\mathcal H$:
\[
\dot F(t)=\{F,\mathcal H\}(x(t),\alpha(t))\qquad \forall\,F\in C^\infty(A^*).
\]
In local coordinates $(x^i,\mu_a)$ this reproduces the standard algebroid Hamilton equations
\[
\dot{x}^i=\rho_a^i(x)\frac{\partial\mathcal H}{\partial \mu_a},\qquad
\dot{\mu}_a=-\rho_a^i(x)\frac{\partial\mathcal H}{\partial x^i}
-C_{ab}^c(x)\mu_c\frac{\partial\mathcal H}{\partial \mu_b}.
\]

\medskip
\noindent\textbf{Step 3 (Hamiltonian system $\Rightarrow$ variational criticality).}
Conversely, let $\alpha(t)$ solve the Hamiltonian system on $A^*$ generated by the Pontryagin Hamiltonian $H$ and define $u(t)=u^*(\alpha(t))$. Then the triple $(x(t),u(t),\alpha(t))$ satisfies the PMP stationarity condition together with the state/adjoint dynamics, that is, the Pontryagin differential equations for the state and costate variables associated with $H$.



%Conversely, let $\alpha(t)$ solve the Hamiltonian system on $A^*$ generated by the Pontryagin Hamiltonian $H$ and define $u(t)=u^*(\alpha(t))$. Then the triple $(x(t),u(t),\alpha(t))$ satisfies the PMP stationarity condition together with the state/adjoint dynamics, namely the coupled Pontryagin equations encoded by the Hamiltonian flow of $H$ on $A^*$.

The PMP conditions are precisely the first-order necessary conditions for criticality of the original constrained variational problem, hence $u(t)$ is a critical control and $x(t)$ a critical trajectory.

\medskip
\noindent\textbf{Step 4 (Poisson invariance and symplectic leaves).}
Since $X_{\mathcal H}$ is a Hamiltonian vector field for the Poisson bracket on $A^*$, its flow preserves the Poisson tensor. Therefore trajectories remain in the integral manifolds of the characteristic distribution, i.e.\ in the symplectic leaves of $A^*$.

\medskip
This establishes the equivalence between the variational and Hamiltonian formulations, and the leafwise confinement of optimal dynamics.
\end{proof}

\end{document}
